After we have generated our ChArUco board we will print it out and take pictures of it. We need to take around 10 pictures and they need to be from different angles. Those pictures need to be put in the \texttt{calibdata/} folder as jpgs. Then \texttt{calib.py} can be run and \texttt{intrinsic.npy} and \texttt{distCoeff.npy} will be generated which are the intrinsic matrix and distortion coefficients also the re-projection error.

\subsubsection{Board parameters and ChArUco detector}
At the top of \texttt{calib.py} we firstly declare the parameters of our board. \texttt{width} and \texttt{height} are the number of rows and culumns that we specified previously. \texttt{square\_len} and \texttt{marker\_len} are the real world sizes of the squares and markers measured in meters. \texttt{dict} is the dictionary number of the ChArUco pattern. With them we can create a \texttt{aruco\_CharucoBoard} object.

\begin{lstlisting}[language=Python]
width = 7
height = 10
square_len = 27.2/1000
marker_len = 14/1000
dict = 12
    
aruco_dict = cv.aruco.getPredefinedDictionary(dict)
board_size = (width, height)
board = cv.aruco.CharucoBoard(board_size, square_len, marker_len, aruco_dict)
\end{lstlisting}

\noindent Next we need a \texttt{aruco\_CharucoDetector} object. For it we need the board object and parameter objects. These are charuco detection, marker and marker refinement parameters. To use refined board we need to set \texttt{tryRefineMarkers} to true in the charuco detection parameters. This way we detect more markers. 

\begin{lstlisting}[language=Python]
refparams = cv.aruco.RefineParameters(minRepDistance = 10.,
									   errorCorrectionRate = 3.,
									   checkAllOrders=True)
    
charucoparms = cv.aruco.CharucoParameters()
charucoparms.tryRefineMarkers = True
    
charuco_detector = cv.aruco.CharucoDetector(board, 
                                            refineParams=refparams,
                                            charucoParams=charucoparms)
\end{lstlisting}


\subsubsection{Read images}
Using \texttt{glob} we get the filenames of images in the \texttt{calibdata/} folder and we loop for each of them. With \texttt{cv.imread()} we read the image then with \texttt{cv.cvtColor(img, cv.COLOR\_BGR2GRAY)} we convert it to grayscale. Simplifying the image can also be done by thresholding, but it gives a larger re-projection error. 

\begin{lstlisting}[language=Python]
images = glob.glob('calibdata/*.jpg')
    
for fname in images:    
    img = cv.imread(fname)
    gray = cv.cvtColor(img, cv.COLOR_BGR2GRAY)
\end{lstlisting}

\subsubsection{Corner detection}
We can detect the corners and markers with the \texttt{detectBoard()} function of the charuco detector object. We append the corner positions to an \texttt{imgpoints} array and we append the found corners to an \texttt{objpoints} array. This array stores where the points are on the board(eg. ((0,0,0), (1,0,0), (2,0,0))). We can also draw the found corners with \texttt{cv.aruco.drawDetectedCornersCharuco(gray, charuco\_corners, charuco\_ids)} and show the result with \texttt{cv.imshow()}.

\begin{lstlisting}[language=Python]
    charuco_corners, charuco_ids, marker_corners, marker_ids = charuco_detector.detectBoard(gray)
    
    imgpoints.append(charuco_corners)

    objp = np.zeros(((height-1)*(width-1),3), np.float32)
    objp[:,:2] = np.mgrid[0:(width-1),0:(height-1)].T.reshape(-1,2)
    
    objtemp = []
    for id in charuco_ids:
        objtemp.append(objp[int(id)])
  
    objpoints.append(np.array(objtemp))

    cv.aruco.drawDetectedCornersCharuco(gray, charuco_corners, charuco_ids)
    cv.imshow('detected corners', gray)
    
\end{lstlisting}

\subsubsection{Camera calibration and saving}
Now we can perform the camera calibration with \texttt{cv.calibrateCamera(objpoints, imgpoints, (w, h), None, None)}. This function returns the intrinsic matrix \texttt{mtx}, distortion coefficients \texttt{dist}, rotational vectors and translation vectors(extrinsic parameters, unique for every image) \texttt{rvec, tvec}. We can save the intrinsic matrix and distortion coefficients for later use using the numpy function \texttt{np.save()}. 

\begin{lstlisting}[language=Python]
ret, mtx, dist, rvecs, tvecs = cv.calibrateCamera(objpoints, imgpoints, (w, h),
                                                    None, None)

np.save('intrinsic.npy', mtx)
np.save('distCoeff.npy', dist)
\end{lstlisting}

\subsubsection{Re-projection error}
With the values that we got from the camera calibration, we can calculate the re-projection error which is a good measure for estimating how good our calibration is. The closer it is to 0 the better. 

\begin{lstlisting}[language=Python]
mean_error = 0
for i in range(len(objpoints)):
    imgpoints2, _ = cv.projectPoints(objpoints[i], rvecs[i], tvecs[i], mtx, dist)
    error = cv.norm(imgpoints[i], imgpoints2, cv.NORM_L2)/len(imgpoints2)
    mean_error += error

print( "Total error: {}".format(mean_error/len(objpoints)) )
\end{lstlisting}